



\begin{ledgroupsized}[r]{120mm}
\footnotesize
\pstart
\noindent\textbf{Überlieferung:}
\pend
\end{ledgroupsized}
\begin{ledgroupsized}[r]{114mm}
\footnotesize
\pstart \parindent -6mm
\makebox[6mm][l]{\textit{L}}%
Konzept: 
LH XXXVII~5 Bl.~74.
Ein Blatt 8\textsuperscript{o}. 
\href{\Scan{LH_37_05_074wz}}%
{Gegenmarke} eines Wasserzeichens in der Blattecke;
Papiererhaltungsmaßnahmen;
Alle Ränder beschnitten.
Zwei Seiten.
\pend
\end{ledgroupsized}
\vspace{5mm}
\begin{ledgroup}
\footnotesize
\pstart
\noindent
\textbf{Datierungsgründe:}
Der Träger von N.\;\ref{60314} enthält die Gegenmarke eines Wasserzeichens. Diese kommt auch im Zeugen \textit{L\textsuperscript{2}} der \glqq Solutio problematis R.\,P. Vannii S.\,J.\grqq\ (N.\;\ref{41823}) vor, den Leibniz während seiner Rückkehr aus Italien abfasste, und zwar zwischen April und Juni 1690 (siehe die Einleitung zu N.\;\ref{intro_vanni-solutio-problematis}). Da Leibniz vor oder während seines Aufenthaltes in Wien (30. April bis Mitte Mai 1690) mit der Arbeit an N.\;\ref{41823} begann, muss ihm zu diesem Zeitpunkt Papier von der Sorte vorhanden sein, die wir auch in N.\;\ref{60314} wiederfinden. Daher teilen beide Stücke einen Terminus post quem (April 1690); ferner erscheint eine Entstehung des kurzen Konzepts N.\;\ref{60314} binnen einer entsprechenden Zeitspanne als wahrscheinlich, was den vorgeschlagenen Terminus ante quem ergibt.
\pend 
\end{ledgroup}
\vspace{8mm}
\pstart\noindent
\href{\Scan{LH_37_05_074r}}%
{\lbrack74~r\textsuperscript{o}\rbrack} 
\pend
\nobreak	%
\begin{edfigure}[0em][0.5em]
\includegraphics[scale=1]{%
gesamttex/edit_VIII-4/images/LH_37_05_074_d1_074r.pdf}
\end{edfigure}
\pstart\noindent
Si nisuum\indexS{nisus} 
\edtext{grave\indexS{grave} sustinentium}{\lemma{}\Bfootnote{grave \bars{impellentium vel \textit{gestr.}} sustinentium \textit{erg.~L}}} 
\textit{AC}, \textit{EC} 
\edtext{(\textls{fig.\;1})}{\lemma{}\Bfootnote{(\textls{fig.\;1}) \textit{erg.~L}}}
concursus\indexS{concursus nisuum grave sustinentium} \textit{C} coincidat cum centro gravitatis\indexS{centrum gravitatis} \textit{G}, 
\edtext{erunt vires sustinentes\indexS{vis sustinens}}{\lemma{erunt}\Bfootnote{%
\textit{(1)}~conatus su %
\textit{(2)}~vires sustinentes~\textit{L}}}
inter se et ad grave\indexS{grave} 
\edtext{ut in}{\lemma{ut}\Bfootnote{%
\textit{(1)}~latera %
\textit{(2)}~in~\textit{L}}}
completo parallelogrammo \textit{CEFA}, 
cujus diagonalis sit verticalis, latera parallelogrammi sunt inter se et ad diagonalem. 
\pend
\pstart 
Quod si 
\edtext{concursus nisuum\indexS{concursus nisuum}}{\lemma{concursus}\Bfootnote{%
\textit{(1)}~gravium sit extra centrum~g %
\textit{(2)}~nisuum~\textit{L}}}
\textit{C} sit extra centrum gravitatis\indexS{centrum gravitatis} (\textit{G})\lbrack,\rbrack\ 
cadat tamen in ejus directionem, 
ut nempe \textit{C}(\textit{G}) sit verticalis\lbrack,\rbrack\ idem est quod 
\edtext{ante, perinde}{\lemma{ante,}\Bfootnote{%
\textit{(1)}~nam %
\textit{(2)}~perinde~\textit{L}}}
enim est ac si pondus reductum in punctum (\textit{G}) suspensum esset ex puncto \textit{C}, 
ubi nihil refert quam longum sit filum.\indexS{filum} 
\pend
\clearpage
\pstart
Sed 
\edtext{et}{\lemma{}\Bfootnote{et \textit{erg.~L}}} 
si concursus nisuum\indexS{concursus nisuum} \textit{C} sit extra lineam directionis\indexS{linea directionis} 
ipsius centri gravitatis\indexS{centrum gravitatis} ((\textit{G})) idem adhuc verum. 
Quod demonstrare cum magnum operae pretium esse videatur 
conficiemus diversis modis idem concludentibus. 
\pend
\begin{edfigure}[1.32em][0.5em]
\includegraphics[scale=1]{%
gesamttex/edit_VIII-4/images/LH_37_05_074_d2_074r.pdf}
\end{edfigure}
\pstart
Incumbat exempli 
\edtext{gratia (\textls{fig.\;2}) circulus}{%
\lemma{gratia}\Bfootnote{%
\textit{(1)}~globus 
\textit{(a)}~ce
\textit{(b)}~duobus plan %
\textit{(2)}~fig.
\textit{(3)}~(\textls{fig.\;2}) circulus~\textit{L}}}
centrum habens \textit{C} duabus 
\edtext{rectis inclinatis\indexS{recta inclinata} \textit{AB}, \textit{ED}; utique}{%
\lemma{rectis}%
\Bfootnote{%
\textit{(1)}~inclinatis; cum
\textit{(2)}~inclinatis \textit{AB}, \textit{ED}; utique~\textit{L}}}
cum nisus sint in normali ad tangentem, concurrent in centro circuli \textit{C}, 
nempe \textit{AC}, \textit{EC}. Quod si jam circuli centrum gravitatis\indexS{centrum gravitatis} sit extra centrum 
\edtext{figurae,\indexS{centrum figurae} quod}{\lemma{figurae,}\Bfootnote{%
\textit{(1)}~q %
\textit{(2)}~imo et ejus directionem, %
\textit{(3)}~quod~\textit{L}}}
utique fieri potest, si constet ex materia heterogenea\lbrack,\rbrack\ poterit ita collocari circulus ut 
\edtext{centrum circuli \textit{C},}{\lemma{centrum}\Bfootnote{%
\textit{(1)}~grav. ejus, \textit{G}, %
\textit{(2)}~circuli \textit{C},~\textit{L}}}
etiam cadat extra lineam directionis\indexS{linea directionis centri gravitatis} ipsius 
centri gravitatis\indexS{centrum gravitatis} \textit{G}\lbrack,\rbrack\
quo facto quaeritur quantum in eo situ premantur plana. 
Ajo eodem prorsus modo premi, ac si centrum gravitatis\indexS{centrum gravitatis}
\edtext{\textit{G}}{\lemma{}\Bfootnote{\textit{G} \textit{erg.~L}}}
esset in centro circuli \textit{C}, seu in puncto concursus\indexS{concursus nisuum}\indexS{punctum concursus nisuum} nisuum.\indexS{nisus} 
Et quidem fatendum 
\edtext{est ob \textit{CA},}{%
\lemma{est}%
\Bfootnote{%
\textit{(1)}~sal
\textit{(2)}~poss
\textit{(3)}~ob
\textit{(a)}~\textit{CE},
\textit{(b)}~\textit{CA},
\textit{(c)}~\textit{CA},~\textit{L}}}
\textit{CE} aequales posse circulum non impedientibus planis sustinentibus\indexS{plana sustinentia} 
(quae perfecte plana seu polita suppono) gyrari circa \textit{C}, donec \textit{G} 
veniat sub \textit{C}, seu donec \textit{C} sit in directione ipsius \textit{G}. Nobis tamen licet disquirere, 
\edtext{quantum}{\lemma{}\Bfootnote{quantum \textit{erg.~L}}}
ante illam conversionem in ipso momento collocationis
\edtext{datae, plana\indexS{planum}}{\lemma{datae,}\Bfootnote{%
\textit{(1)}~extra %
\textit{(2)}~plana~\textit{L}}}
sustentent. Ut igitur ostendam nihil interesse, 
\edtext{utrum \textit{C}}{\lemma{utrum}\Bfootnote{%
\textit{(1)}~\textit{G} sit extra %
\textit{(a)}~\textit{C} an non %
\textit{(b)}~dire %
\textit{(2)}~et %
\textit{(3)}~\textit{C}~\textit{L}}}
sit extra directionem ipsius \textit{G} an non; considerandum est motum\indexS{motus compositus ex duobus} 
ipsius \textit{G} posse intelligi compositum ex duobus, uno 
\edtext{parallelo et}{%
\lemma{parallelo}%
\Bfootnote{%
\textit{(1)}~ipsi \textit{C}, altero
\textit{(2)}~et~\textit{L}}}
aequali ipsi \textit{C}, altero circulari\indexS{motus circularis} circa \textit{C}, ex his 
\href{\Scan{LH_37_05_074v}}%
{\lbrack74~v\textsuperscript{o}\rbrack}  
posterior nihil mutat in planis, quia salvo eorum situ fieri potest. Itaque prior solus in rationes veniet, 
ac proinde perinde est ac si \textit{G} caderet in \textit{C}, cum idem sit 
\edtext{ipsius \textit{G} motus}{\lemma{ipsius \textit{G}}\Bfootnote{%
\textit{(1)}~descensus seu %
\textit{(2)}~motus~\textit{L}}}
ad rem faciens qui ipsius \textit{C}. 
\pend 
\begin{edfigure}[1em][0.5em]
\includegraphics[scale=1]{%
gesamttex/edit_VIII-4/images/LH_37_05_074_d3_074v.pdf}
\end{edfigure}
\pstart 
Jam quod de globo\indexS{globus} vel circulo demonstravimus, 
ubi rectae \textit{CA}, \textit{CE} sunt aequales, id de aliis quibuscunque corporibus 
quae duobus planis incumbunt,\indexS{corpus duobus planis incumbens} demonstrari potest. 
Hoc modo remotis caeteris 
\edtext{\textls{fig.\;3}}{\lemma{}\Bfootnote{\textls{fig.\;3} \textit{erg.~L}}} 
considerentur punctum concursus\indexS{punctum concursus nisuum} \textit{C}, 
et rectae nisuum\indexS{nisus} \textit{CA}, 
\edtext{\textit{CF}, perpendiculares}{%
\lemma{\textit{CF},}%
\Bfootnote{%
\textit{(1)}~perinde ac si
\textit{(2)}~perpendiculares~\textit{L}}}
ad plana\indexS{planum inclinatum} \textit{AB}, \textit{FG}; perinde ac si eae duae rectae essent rigidae. 
Jam si \textit{CA}, \textit{CF} sint inaequales sumatur alia \textit{CE}, pars ipsius \textit{CF}, 
innitens plano \textit{ED} parallelo ad planum \textit{FG}, 
patet nihil referre quanta sit longitudo rectae \textit{CE}, vel \textit{CF}, 
seu eandem manere pressionem plani\indexS{pressio plani} \textit{AB}, et eam quae est plani \textit{ED} 
esse etiam pressionem plani\indexS{pressio plani} \textit{FG}. 
Itaque cum in casu plani \textit{CE} locum habeat proportio quam initio stabilivimus\lbrack,\rbrack\ 
habebit eadem quoque locum in casu inaequalium rectarum concurrentium, nec referet
utrum \textit{CAE}, sit circulus an Ellipsis vel alia figura quaevis. Et in genere nihil referet 
\edtext{\lbrack utrum\rbrack}{\lemma{}\Bfootnote{utrum \textit{erg.~Hrsg}}}
centrum gravitatis\indexS{centrum gravitatis} coincidat 
in punctum concursus\indexS{punctum concursus nisuum} nisuum\indexS{nisus} an non. 
\edtext{Si plura}{\lemma{Si}\Bfootnote{%
\textit{(1)}~plures %
\textit{(2)}~plura~\textit{L}}}
sint plana duobus remedium mox afferemus. 
\pend
\pstart
Aliter idem 
\edtext{sic probatur}{\lemma{sic}\Bfootnote{%
\textit{(1)}~demonstratur %
\textit{(2)}~probatur~\textit{L}}}
etsi minore nonnihil firmitate\indexS{firmitas} planum \textit{AB} renisu\indexS{renisus} suo 
\edtext{(quasi arcu tenso),}{\lemma{}\Bfootnote{(quasi arcu tenso) \textit{erg.~L}}}
conatur
\edtext{impellere corpus \textit{CAD}, ita ut quodlibet ejus punctum moveatur}{\lemma{impellere}\Bfootnote{%
\textit{(1)}~\textit{C} ut recta %
\textit{(2)}~corpus \textit{CAD}, \lbrack...\rbrack\ moveatur~\textit{L}}}
secundum rectam \textit{AC} et parallelas, perinde ac si 
exiguum latus polygoni quo incumbit \textit{AB} impelleretur 
ita ut quodlibet ejus punctum moveretur parallele ad \textit{AC}. Eodem 
\edtext{modo planum}{\lemma{modo}\Bfootnote{%
\textit{(1)}~p %
\textit{(2)}~quodlibet punctum corporis \textit{C} %
\textit{(3)}~planum~\textit{L}}}
\textit{ED} renisu\indexS{renisus} 
\edtext{suo conatur}{\lemma{suo}\Bfootnote{%
\textit{(1)}~(quasi %
\textit{(2)}~conatur~\textit{L}}}
impellere corpus \textit{CAD} ita ut quodlibet punctum moveatur secundum 
\edtext{rectam \textit{EC}}{\lemma{rectam}\Bfootnote{%
\textit{(1)}~\textit{AC} %
\textit{(2)}~\textit{EC}~\textit{L}}}
et parallelas. Et proinde punctum \textit{G}, quod est centrum gravitatis corporis\indexS{centrum gravitatis corporis} 
et toti corpori aequivalet, recipiet duos nisus\indexS{nisus}\lbrack,\rbrack\ unum in parallela ipsi \textit{AC}, 
alterum in parallela ipsi \textit{EC}, quorum compositio\indexS{compositio nisuum} dabit nisum ipsius \textit{G} verticalem, 
seu aequivalebit ejus gravitati ut in parallelogrammo figurae 1. 
Potest tamen aliquid objici quod facit ut huic argumento satis per se fidi non posset. Nam 
\edtext{(nisi \textit{ED} sit normalis ad \textit{AB})}{\lemma{}\Bfootnote{(nisi \textit{ED} sit normalis ad \textit{AB}) \textit{erg.~L}}} 
ipsum planum\indexS{planum sustinens} \textit{ED} obstat impulsui plani \textit{AB} et viceversa, quod in aliis conatuum compositionibus\indexS{compositio conatuum} non fit. Praeterea 
\edtext{grave \textit{C}}{\lemma{grave}\Bfootnote{%
\textit{(1)}~impulsum conatu ecc %
\textit{(2)}~\textit{C}~\textit{L}}}
impulsum 
\edtext{solo}{\lemma{}\Bfootnote{solo \textit{erg.~L}}}
nisu Eccentrico\indexS{nisus eccentricus}
\textit{AC}, revera non moveretur in parallelis ipsi \textit{AC}, quia simul \textit{G} 
conabitur circulari circa \textit{A} radio \textit{AG}. 
Cui licet nunc obstet alterum planum, videbitur tamen hic 
\edtext{ipse ipsius}{\lemma{ipse}\Bfootnote{%
\textit{(1)}~nisus %
\textit{(2)}~ipsius~\textit{L}}}
\textit{G} circa \textit{A} circulare conantis nisus\indexS{nisus} in planum\indexS{planum sustinens} alterum \textit{ED} 
debere venire in rationes. Idemque est de impulsu ipsius \textit{EC}, 
nisi \textit{G} cadat in rectam \textit{EC}.
\pend
\begin{edfigure}[1em]%
\includegraphics[scale=1]{%
gesamttex/edit_VIII-4/images/LH_37_05_074_d4_074v.pdf}
\edcaption[-4.5ex]{%
\edtext{}{%
\lemma{\textls{fig.\;4}}%
\Cfootnote[nonum]{Ein gestrichener Entwurf zum Diagramm wird nicht wiedergegeben.}}}
\end{edfigure}
\pstart 
\edtext{Sed et alia demonstratio melior suppetit.}{\lemma{Sed}\Bfootnote{\bars{et \textit{erg.}} alia %
\textit{(1)}~adhuc superest %
\textit{(2)}~demonstratio \bars{melior \textit{erg.}} suppetit.~\textit{L}}}
Fig.\;4 ajo primum, eodem modo premi planum \textit{AB}, 
sive corpus \textit{EBA} incumbat in puncto \textit{A}, sive in puncto \textit{B}, sive in ambobus. 
Ducta enim \textit{E3G2} in centrum gravitatis\indexS{centrum gravitatis} \textit{G}, 
\edtext{et ad \textit{AB}}{\lemma{et}\Bfootnote{%
\textit{(1)}~id pat %
\textit{(2)}~ad \textit{AB}~\textit{L}}}
normalibus \textit{A2}, \textit{B3}, patet \textit{EG} sive quatenus fluit
per \textit{ED} circa \textit{E} eodem modo agere in \textit{A}, quo  
\edtext{in \textit{B}\lbrack,\rbrack\ nam etsi \textit{B} videatur facilius moveri 
tamen tanto minus promovetur\lbrack,\rbrack\ unde}{\lemma{}\Bfootnote{%
in \textit{B}
\bars{%
\textit{(1)}~non enim %
\textit{(2)}~nam etsi \lbrack...\rbrack\ minus promovetur \textit{erg.}}
unde~\textit{L}}} 
si ambobus simul sive quatenus volvitur dimidia tantum vis 
est in uno puncto tota in integro plano\indexS{planum integrum} 
v\angles{olve}nte, non ergo refert. 
\pend
\pstart
Hinc jam 
\edtext{sequitur, \textls{quotcunque}}{\lemma{sequitur,}\Bfootnote{%
\textit{(1)}~nisu utcun %
\textit{(2)}~\textls{quotcunque}~\textit{L}}}
\textls{assumtis planis}, et nisibus utcunque eccentricis,\indexS{nisus eccentricus} 
sufficere ut ex centro 
\edtext{gravitatis\indexS{centrum gravitatis} agantur}{%
\lemma{gravitatis}%
\Bfootnote{%
\textit{(1)}~agat
\textit{(2)}~agantur~\textit{L}}}
normales in plana, et perinde fore, ac si ubi normales 
ex centro gravitatis\indexS{centrum gravitatis} planis\indexS{planum sustinens} occurrent, esset contactus.
\pend
